\subsection{\textcolor{Salmon}{Beispielrechnung zur Hypothesen-Auswertung anhand von t-Tests}}
\label{ch:Ergebnisse:sec:Hypothesen_Auswertung:sub:Bsp_TTest}
\textcolor{blue}{\textit{einleitende Worte}}
% Wikipedia:
%Der t-Test ist ein Begriff aus der mathematischen Statistik, er bezeichnet eine Gruppe von Hypothesentests mit t-verteilter Testprüfgröße. Oft ist jedoch mit dem t-Test der Einstichproben- bzw. Zweistichproben-t-Test auf einen Mittelwertunterschied gemeint. 

\noindent
\textcolor{cyan}{Was macht der T-Test?}\\
- Differenzen bilden\\
- davon Mittelwert und Standardfehler berechnen\\
- Aussage treffen, wie wahrscheinlich es ist, dass dieses Ergebnis zufällig entstanden ist\\
\textcolor{red}{- TODO: einseitig / zweiseitig definieren}

\begin{quote}
    t-Test. Verfahren zur Überprüfung des Unterschiedes zweier Stichprobenmittelwerte. Man unterscheidet den \\
    -> t-Test für unabhängige Stichproben (Vergleich der Mittelwerte von Stichproben aus zwei verschiedenen Populationen) und den\\
    -> t-Test für abhängige Stichproben (Vergleich zweier Mittelwerte einer Variablen, die an derselben Stichprobe zu zwei verschiedenen Zeitpunkten oder an »Matched Samples« erhoben wurden).
    - \cite{Statistik:Boertz:2006}
\end{quote}

\noindent
\textcolor{cyan}{- weitere Quellen: \cite{Statistik:Boertz:2010} \cite{Statistik:Sedlmeier:2024}}\\

\noindent
\begin{table}[H]
Das nachfolgende Beispiel nutzt die Daten von Frage 1 der Nutzerbefragung (s. \autoref{fig:pruemper_all}). \\
\begin{tabular}{|l|c|c|c|c|c|c|c|c|c|c|c|}
    \hline
    System      & \multicolumn{11}{l|}{Antworten der Nutzer} \\ \hline
    alt         &  2 & 2 & 1 & -1 & 1 & -2 & 2 & -3 & -2 & -1 & -1 \\ \hline
    neu         &  0 & 3 & 1 &  2 & 2 &  1 & 3 &  3 &  3 &  1 & -2 \\ \hline
    \hline
    Differenz x  & -2 & 1 & 0 &  3 & 1 &  3 & 1 &  6 &  5 &  2 & -1 \\
    \hline
\end{tabular}
\end{table}

\noindent
Differenzen mit Ergebnis = 0 werden nicht betrachtet \textcolor{red}{TODO: Quelle!}\\
Hypothesen:\\
\textcolor{red}{
$H0:= alt \geq neu$\\
$H1:= alt < neu$\\
}

\noindent
Es wird ein einseitiger, gepaarter t-Test durchgeführt. Die Vorraussetzungen dafür sind folgende:
\begin{itemize}
    \item es liegen zwei abhängige Stichproben vor $\rightarrow$ jede Person beantwortet alle Fragen für zwei unterschiedliche Systeme
    \item die Daten sind metrisch
    \item Differenzen sind normalverteilt $\rightarrow$ wurde überprüft mithilfe des Online-Rechners von StatistikGuru \cite{normalverteilung_rechner}
\end{itemize}

\noindent
Anzahl der Fälle:

$\displaystyle n = 10$ \\

\noindent
Freiheitsgrade:

$\displaystyle d_f = n-1 = 9$  \\

\noindent
Mittelwert:
\begin{flalign*}
\indent
\displaystyle
\Bar{x}   &= \frac{\sum{x_{diff}}}{n} &\\
          &= \underline{\underline{1,9}}
\end{flalign*}

%\newpage
\noindent
Standardabweichung:
\begin{flalign*}
\indent
\displaystyle
    s %&= \sqrt{ \frac{ \Xsum\limits_{i=0}^n (x_n - \Bar{x})^2  }{ n-1 } } &\\
    &= 
    \sqrt
    {
        \frac
        {
            \Xsum\limits_{x=0}^n (x_n - \Bar{x})^2
            % \Xsum_{x=0}^n (x_n - \Bar{x})^2
        }
        {
            n-1
        }
    }&\\
%    =%&=
%    \sqrt
%    {
%        \frac
%        {
%            (x_1 - \Bar{x})^2 +
%            (x_2 - \Bar{x})^2 +
%            (x_3 - \Bar{x})^2 +
%            ... +
%            (x_n - \Bar{x})^2
%        }
%        {
%            n-1
%        }
%    } &\\
    &=
    \sqrt
    {
        \dfrac
        {
            \splitdfrac{
                (-2 - 1,9)^2
                + ( 1 - 1,9)^2
                + ( 3 - 1,9)^2
                + ( 1 - 1,9)^2
                + ( 3 - 1,9)^2
            }{
                + ( 1 - 1,9)^2
                + ( 6 - 1,9)^2
                + ( 5 - 1,9)^2
                + ( 2 - 1,9)^2
                + (-1 - 1,9)^2
            }
        }
        {
            10 - 1
        }
    } &\\
    &=
    \sqrt
    {
        \dfrac
        {
            \splitdfrac{
                (-3,9)^2
                + (-0,9)^2
                + 1,1^2
                + (-0,9)^2
                + 1,1^2
            }{
                + (-0,9)^2
                + 4,1^2
                + 3,1^2
                + 0,1^2
                + (-2,9)^2
            }
        }
        {
            10 - 1
        }
    } &\\
    &=
    \sqrt
    {
        \dfrac
        {
            \splitdfrac{
                15,21
                + 0,81
                + 1,21
                + 0,81
                + 1,21
            }{
                + 0,81
                + 16,81
                + 9,61
                + 0,01
                + 8,41
            }
        }
        {
            9
        }
    } &\\
    &= \sqrt{\frac{54,9}{3}} = \sqrt{6,1} \approx \underline{\underline{2,4698}}
\end{flalign*}

\noindent
Standardfehler des Mittelwerts:\\

$ \displaystyle
s_e = \frac{s}{\sqrt{n}}
    = \frac{2,4698}{\sqrt{10}}
    \approx \underline{\underline{0,7810}}
$\\

\noindent
Teststatistik:\\

$ \displaystyle
t  = \frac{\Bar{x} - 0}{s_e}
   = \frac{\Bar{x}}{s_e}
   = \frac{1,9}{0,781}
   \approx \underline{\underline{2,432778}}
$\\

\noindent
Wahrscheinlichkeit: \textcolor{red}{TODO: Verweis auf die Statistik-Tabelle}\\

$ \displaystyle
p \approx \underline{\underline{0,019471718}}
$ (aus Excel)\\

\noindent
Der t-Test kann mithilfe von Microsoft Excel durchgeführt werden. Die Formel für einen einseitigen, gepaarten t-Test ist:
$$=TTEST(Matrix_{alt}, Matrix_{neu},1,1)$$
Alle weiteren t-Test Ergebnisse wurden über diesen Weg erreicht und sind in \autoref{tab:usability_test:Differenz+ttest_Fehler} für die Fehlerquotenvergleich des Usability-Tests, in \autoref{tab:usability_test:Differenz+ttest_Zeit} für den Zeitvergleich des Usability-Tests und in \autoref{tab:Nutzerbefragung:Differenz+ttest} für den Vergleich der Ergebnisse der Nutzerbefragung aufgeführt.
